% CCSS 7.NS.A.1.B — Adding Integers
\section{Adding Integers}

\begin{learningGoals}
\begin{itemize}
    \item Add integers using number line reasoning
    \item Understand that a number plus its opposite equals zero
\end{itemize}
\end{learningGoals}

\begin{conceptBox}{Rules for Adding Integers}
\begin{itemize}[leftmargin=*]
    \item \textbf{Same signs:} Add the absolute values. Keep the sign.
    
    $5 + 3 = 8$ \qquad $(-5) + (-3) = -8$
    \item \textbf{Different signs:} Subtract the smaller absolute value from the larger. Keep the sign of the larger.
    
    $7 + (-4) = 3$ \qquad $(-7) + 4 = -3$
    \item \textbf{Additive inverses:} A number plus its opposite is always $0$.
    
    $9 + (-9) = 0$
\end{itemize}
\end{conceptBox}

\begin{workedExample}{Adding on a Number Line}
Find $(-3) + 5$.

Start at $-3$. Adding $5$ means move $5$ units \textbf{right}.

$-3 \to -2 \to -1 \to 0 \to 1 \to 2$

\answerHighlight{$(-3) + 5 = 2$}
\end{workedExample}

\tipBox{Think of positive numbers as steps right and negative numbers as steps left. The number line never lies!}

\begin{practiceBox}{Adding Integers}
\resetProblems

\prob $(-8) + (-5)$
\answerExplain{$-13$}{Both are negative, so add the absolute values: $8 + 5 = 13$. Keep the negative sign: $-13$.}

\prob $(-12) + 7$
\answerExplain{$-5$}{Different signs, so subtract: $12 - 7 = 5$. The larger absolute value is $12$ (negative), so the answer is $-5$.}

\prob $15 + (-15)$
\answerExplain{$0$}{These are opposites (additive inverses). A number plus its opposite always equals $0$.}

\prob The temperature is $-6°$C in the morning. By noon it rises $14°$. What is the noon temperature?
\answerExplain{$8°$C}{Start at $-6$ and add $14$: $(-6) + 14 = 8$. Different signs, so subtract $14 - 6 = 8$, and the larger absolute value is positive.}

\prob $(-9) + (-4) + 6$
\answerExplain{$-7$}{First add the negatives: $(-9) + (-4) = -13$. Then add $6$: $(-13) + 6 = -7$. Subtract $13 - 6 = 7$, keep the negative sign.}

\end{practiceBox}
